%
% (c) Rafael Rodríguez Galván, Guillem Borrell
% Se autoriza la utilización de este documento según la licencia GFDL,
% sin secciones invariantes, texto de portada ni de respaldo.
%

\documentclass[12pt,xcolor=svgnames,spanish]{beamer}
\usepackage[T1]{fontenc}
%\usepackage{lmodern}
%\usepackage[scaled]{helvet}
\usepackage{babel}
\usepackage[utf8]{inputenc}
\usepackage{tabularx}
\usepackage{longtable}
\usepackage{amsmath}
\usepackage{smbslides}

\mode<presentation>
{
%  \usetheme{Frankfurt}
  \usetheme{Madrid}
  \usecolortheme{UCA}

%  \useoutertheme{infolines}
  %\setbeamersize{text margin left=1cm,text margin right=1cm}
 \useinnertheme{circles}
  %\usecolortheme{wolverine}

 \addtobeamertemplate{block begin}{}
  %{\setbeamercolor{local structure}{fg=block title.bg}}

  % Parece no tener efecto cuando se presenta el índice. ¿Es por el modo?
  %\setbeamercovered{dynamic}
}

\AtBeginSection[] % Do nothing for \section*
{
  \begin{frame}<beamer>
   \frametitle{Sección \thesection. \insertsection}
    \tableofcontents[currentsection]
  \end{frame}
}

\AtBeginSubsection[] % Do nothing for \subsection*
{
  \begin{frame}<beamer>
    \frametitle{Sección \thesection.\thesubsection 
      \par \small{\insertsection: \insertsubsection}}
    \tableofcontents[currentsection,currentsubsection]
  \end{frame}
}

\title{SM: Cálculo numérico con Octave}

\pgfdeclareimage[height=1.76mm]{ccsa}{img/cc-bysa}
\logo{\pgfuseimage{ccsa}}
\author{J. Rafael Rodríguez Galván}
\institute[OSLUCA]{Departamento de Matemáticas. Universidad de Cádiz)}
\date[i-MATH, Santiago 2010]{Sofware en Matemáticas. Máster Interuniversitario en Matemáticas}


% \includeonly{toma_contacto}
% \includeonly{matrices}
% \includeonly{programacion}
\includeonly{toma_contacto, matrices, programacion, graficos}

%%%%%%%%%%%%%%%%%%%%
\begin{document}
%%%%%%%%%%%%%%%%%%%%

\maketitle

\begin{frame}
\vfill\tiny
(c) Rafael Rodríguez Galván, Guillém Borrell.
Se autoriza la utilización de este documento según la licencia GFDL,
sin secciones invariantes, texto de portada ni de respaldo.
\par
\hfill\url{http://ciencialibre.forja.rediris.es}
\end{frame}

%___________________________________________________________________________

\section{Toma de contacto con Octave}

\subsection{Introducción a Octave}

% ___________________________________________________________________________

\hypertarget{qu-es-octave-i}{}
\pdfbookmark[0]{Qué es Octave}{qu-es-octave-i}
\begin{frame}\frametitle{Qué es Octave}
\begin{itemize}
\item {} 
  \textbf{Lenguaje} interpretado de alto nivel, orientado al \textbf{cálculo matricial}

\item {} 
  \textbf{Intérprete} para este lenguaje

\item {} 
Licencia libre (GPL)
\begin{itemize}
\item {} 
Concebido en 1988 por John W. Eaton

\item {} 
Amplia comunidad de usuarios (\href{http://www.octave.org}{http://www.octave.org})

\item {} 
Código fuente C++
\begin{itemize}
\item {} 
Potente biblioteca de cálculo matricial (GPL)

\item {} 
Fácilmente extensible en C++

\end{itemize}

\item {} 
Disponible para numerosos sistemas y arquitecturas

\end{itemize}


\item Octave \emph{\bf no} está especializado en \emph{\bf cálculo simbólico}
    \begin{itemize}
    \item {} Para ello, existen numerosas herramientas libres como
       \textbf{Python}+\textbf{SymPy}, \textbf{Maxima}... 
    \end{itemize}
    
\item 
    Existen muchos otros entornos de cálculo numérico libres:
    \begin{itemize}
      \item {} \textbf{Python}+\textbf{NumPy}/\textbf{Scipy}, \textbf{Julia}, Scilab, Euler, Freemat...    
    \end{itemize}
    
\end{itemize}

\end{frame}


% ___________________________________________________________________________

\hypertarget{potencialidad-de-octave}{}
\pdfbookmark[0]{Potencialidad de Octave}{potencialidad-de-octave}
\begin{frame}\frametitle{Potencialidad de Octave}
\begin{itemize}
\item {} 
Curva de acceso poco empinada

\item {} 
Incluye ``de serie'':
\begin{itemize}
\item {} 
Operaciones matriciales

\item {} 
Resolución de (sistemas de) ecuaciones lineales y no lineales

\item {} 
Cálculo de autovalores y autovectores

\item {} 
Cálculo de raíces de ecuaciones no lineales

\item {} 
Aproximación numérica de funciones y datos (interpolación,
FFT, funciones spline,...)

\item {} 
Cálculo numérico de derivadas e integrales

\item {} 
Resolución numérica de ecuaciones diferenciales

\item {} 
Representación de gráficas 2D y 3D 

\item {} 
...

\end{itemize}

\item Compatibilidad con Matlab...

\end{itemize}

\end{frame}

% ___________________________________________________________________________


\hypertarget{octave-y-matlab}{}
\pdfbookmark[0]{Octave y Matlab}{octave-y-el-lenguaje-matlab}
\begin{frame}\frametitle{Octave y Matlab}
  \begin{itemize}
  \item 
    Alta compatibilidad con el ``\textbf{lenguaje Matlab}''
    \begin{itemize}
    \item {} ventajas
      \begin{itemize}
      \item fácil migración desde Matlab
      \item hereda puntos fuertes
      \end{itemize}
    \item {} inconvenientes
      \begin{itemize}
      \item freno a la innovación propia
      \item hereda puntos débiles
      \end{itemize}
    \end{itemize}
  
  \item ¿Dónde se pierde la compatibilidad?
    \begin{itemize}
    \item Extensiones al lenguaje Matlab
    \item ``Toolkits''
    \end{itemize}

  \item Filosofía propia: ``\emph{lenguaje de scripting científico}''
    \begin{itemize}
    \item Se complementa con otras herramientas libres
    \item Entorno gráfico similar a Matlab + integración con otros editores/entornos... 
    \end{itemize}

  \end{itemize}
\end{frame}

% % ___________________________________________________________________________


\hypertarget{qu-es-octave-ii}{}
\pdfbookmark[0]{Octave y Matlab}{qu-es-octave-ii}
\begin{frame}\frametitle{Lenguajes interpretados en cálculo científico}

\begin{itemize}
\item {} 
Lenguajes \emph{interpretados}:
\begin{itemize}
\item {} 
Reducción del ciclo de desarrollo (relación esfuerzo/resultado)

\item {} 
Polivalencia, acceso a herramientas de alto nivel

\item {}
Concentrarse en las dificultades del problema (no en
  las del lenguaje)

\end{itemize}

\item {} 
Lenguajes \emph{compilados}: Potencia de cálculo

\item {} 
Lenguajes interpretado como lenguajes ``\emph{pegamento}'': lo mejor de dos mundos
\begin{itemize}
\item 90\% del código: lenguaje interpertado
\item 10\% restante (zonas críticas en potencia de calculo): lenguaje
  compilado
\end{itemize}
\item Para tareas científicas:
\begin{itemize}
\item \textbf{Octave}
\item Python
\item ...
\end{itemize}
\end{itemize}

\end{frame}

% % ___________________________________________________________________________

% \hypertarget{interfaces}{}
% \pdfbookmark[0]{Interfaces}{interfaces}
% \begin{frame}\frametitle{Interfaces}
% \begin{itemize}
% \item {} 
% Octave incluye una interfaz muy sencilla en forma de
% consola (no gráfica)

% \item {} 
% Existen algunas alternativas:
% \begin{itemize}
% \item {} 
% Modo Octave \textbf{Emacs}

% \item {} 
% \textbf{qtOctave}

% \item {} 
% Otras

% \end{itemize}

% \end{itemize}

% \end{frame}


\subsection{Primeros pasos}

\hypertarget{primera-sesi-n-con-qtoctave}{}
\pdfbookmark[0]{Primera sesión con Octave}{primera-sesi-n-con-qtoctave}
\begin{frame}
\frametitle{Primera sesión con Octave}
Arrancamos Octave en modo gráfico...
\begin{itemize}

\item {} 
Distintas secciones

\begin{itemize}

\item {} 
  Explorador de archivos

\item {} Espacio de trabajo: edición de variables
\item {} Historial de comandos
\item {} 
 {Ventana de ``comandos''}: usando Octave como una calculadora

\end{itemize}

\item Otras utilidades:
\begin{itemize}
\item Editor integrado
  \item {Acceso a la ayuda de Octave}%
  \footnote{\tiny\url{chttps://docs.octave.org/interpreter/index.html}}
\item Entorno mejorado para la representación de gráficos
\item ...
\end{itemize}

\end{itemize}

\end{frame}



% ___________________________________________________________________________

\begin{frame}[containsverbatim]
\frametitle{Octave como una calculadora (I)}

\begin{itemize}
\item Operadores numéricos habituales: 
  \par \verb|+|, \verb|-|, \verb|*|, \verb|/|, \verb|^|
\item Paréntesis para modificar la preferencia
\item El resultado se almacena en la variable \verb|ans|
\item Re-editar órdenes: $\uparrow$, $\downarrow$, \emph{(Mayúsculas
    +)} $\rightarrow$, $\leftarrow$ 
\end{itemize}

\small
\begin{Octave}
>> 1+3*1.5
ans =  5.5000
>> (2+3)^10
ans =  9765625
>> ans/5^7
ans =  125
\end{Octave}


\end{frame}


\begin{frame}[containsverbatim]
\frametitle{Octave como una calculadora (II)}
Estilo Matlab...
\begin{itemize}
\item {} 
Comillas simples para introducir una cadena de texto

\item {} 
Porcentaje para comentarios

\item {} 
Tres puntos para continuación de línea

\item {} 
Punto y coma para un retorno de carro sin salida

\end{itemize}

\small
\begin{Octave}
>> 'Esto es una cadena de texto'
ans = Esto es una cadena de texto
>> 20000000 + 3100020500 +...
> 13
ans =  3.1200e+09
>> 2+2; % No se muestra el resultado
\end{Octave}

\end{frame}

\begin{frame}[containsverbatim]
\frametitle{Octave como una calculadora (II)}
Estilo Octave (no compatible!)...
\begin{itemize}
\item {} 
Comillas dobles para introducir una cadena de texto

\item {} 
Almohadilla (\verb|#|) para comentarios

\item {} 
Barra invertida (\verb|\|) para continuación de línea
\end{itemize}

Similar a otros lenguajes de scripting: python, bash...

\small
\begin{Octave}
>> ``Esto es una cadena de texto''
ans = Esto es una cadena de texto
>> 20000000 + 3100020500 +\
> 13
ans =  3.1200e+09
>> 2+2; # No se muestra el resultado
\end{Octave}

\end{frame}


\begin{frame}[containsverbatim]
\frametitle{Más madera...}
\small
\begin{Octave}
>> 1/7          % Aritmetica punto fijo
ans =  0.14286
>> format long  % Mostrar doble precisi'on
>> 1/7 
ans =  0.142857142857143
>> format       % O bien 'format short'
>> sin(pi)      %  sen(pi) = 0 ?
ans =  1.2246e-16
>> eps          % El "cero" de la maquina
ans =  2.2204e-16
>> x=1+2*i; y=2+3*i;  % Asignacion, complejos
>>  x*y
ans = -4 + 7i
\end{Octave}
\end{frame}

\begin{frame}[containsverbatim]
\frametitle{Variables}
Órdenes:
\begin{itemize}
\item \verb|who|: listado de variables
\item \verb|whos|: listado detallado
\item \verb|clear x y ...|: borrar variables \verb|x|, \verb|y|,...
\item \verb|clear all|: borrarlas todas
\item \verb|help orden|: información sobre una orden de octave
\item \verb|save fichero x y ...|\footnote{Octave puede grabar en numerosos
    formatos. Véase la ayuda}
\item \verb|load fichero x y ... | ó \verb|load fichero|
\end{itemize}
Variables reservadas:
\begin{itemize}
\item \verb|ans|, \verb|eps|, \verb|inf|, \verb|nan|, \verb|pi|,
  \verb|i|, \verb|j|, \verb|nargin|, \verb|nargout|, \verb|realmin|,
  \verb|realmax|, \verb|flintmax|, \verb|varargin|, \verb|varargout|
\end{itemize}
\end{frame}

\begin{frame}[containsverbatim]
\frametitle{Ejercicios}
\begin{enumerate}
\item
  Definir las variables:
  \begin{itemize}
  \item $\displaystyle x=e^{\pi-1}$
  \item $\displaystyle y=\frac{1}{1+\frac{1}{x+2}}$
  \item $\displaystyle z=\sqrt{x+y}$
  \item $\displaystyle t=\log(x+y+z)$
  \end{itemize}
\item Almacenarlas en un fichero (binario) llamado ``variables.dat''
\item Borrar todas las variables y volverlas a cargar desde el fichero
\end{enumerate}
\bigskip
``Para subir nota'': Grabar estas variables en formato texto (\texttt{save -ascii}) llamado ``variables.txt'':
\end{frame}

% ___________________________________________________________________________

%%% Local Variables: 
%%% mode: latex
%%% TeX-master: "modulo_octave"
%%% End: 


\section{Matrices y álgebra matricial}
% ___________________________________________________________________________

\hypertarget{secuencias}{}
\pdfbookmark[0]{Secuencias}{secuencias}

\begin{frame}[containsverbatim]
\frametitle{Secuencias}
\small
\begin{itemize}
\item {} 
Son argumentos dedicados a contar.

\item {} 
Uso:
\begin{itemize}
\item contador en bucles
\item definir rangos de valores
\item crear vectores fila.
\end{itemize}

\end{itemize}

Ejemplo:
\begin{Octave}
  >> 1:5 
  ans =
    1 2 3 4 5
  >> 0:2:10
  ans =
    0 2 4 6 8 10
  >> 0:-3:-10
  ans =
    0 -3 -6 -9
\end{Octave}

\end{frame}


\begin{frame}[containsverbatim]
\frametitle{Nuestro primer script desde qtOctave}

\begin{itemize}

\item Editar y grabar en disco un fichero llamado
    ``tablavalores.m'' con el siguiente contenido:

\begin{Octave}[morekeywords={x,y}]
  x=-1:0.1:1; 
  y=1./(1+x.*x);
\end{Octave}

\item Los operadores \verb|.*| y \verb|./| multiplican y dividen
  elemento a elemento.

\item {} En la ventana de comandos, ejecutar el script y pedir los valores
  de \emph{x} e \emph{y}

\begin{Octave}[morekeywords={tablavalores,x,y}]
  >> tablavalores 
  >> x 
    (...valores de x...)  
  >> y
    (...valores de y...)
\end{Octave}


\item Otras funciones disponibles disponibles en la ventana de comandos:
  \texttt{ls}, \texttt{pwd}, \texttt{cd}
\end{itemize}





\end{frame}


% % ___________________________________________________________________________

%\subsection{Matrices}
% ___________________________________________________________________________

\hypertarget{tipos-num-ricos}{}
\pdfbookmark[0]{Tipos numéricos}{tipos-num-ricos}
\begin{frame}[containsverbatim]
\frametitle{Tipos numéricos}

\begin{itemize}
\item 
  \emph{Por omisión}, Octave (al igual que Matlab) emplea un único
  tipo numérico: 
\begin{center}
  \Large
  \textbf{\alert{reales de doble precisión}}
  \\[0.5em]
  \normalsize
  (arrays n-dimensionales de)
  \\[0.4em]~
\end{center}


\item Cualquier otro tipo tendrá que asignarse explícitamente:
\begin{flushright}
  \verb|int8|, \verb|int16|, \verb|int32|, \verb|int64|
  \\
  \small \verb|uint8|, \verb|uint16|, \verb|uint32|, \verb|uint64|
  \\
  \verb|single|, \verb|double|
\end{flushright}

\item Ejemplo: \verb|i=int8(32)|
\end{itemize}
\end{frame}



% ___________________________________________________________________________

\hypertarget{para-diferenciar-filas-y-columnas}{}
\pdfbookmark[1]{Para diferenciar filas y columnas}{para-diferenciar-filas-y-columnas}
\begin{frame}[containsverbatim]
\frametitle{Para escribir matrices}

\textbf{Notación:}
\begin{itemize}
\item Se utilizan corchetes
\item El espacio en blanco o la coma separan columnas en una fila
\item El retorno de carro o el punto y coma separan filas
\end{itemize}

\textbf{Ejemplo} 
\begin{Octave}
>> A=[1,2; 3,4]
A =
   1   2
   3   4
\end{Octave}

\textbf{Ejercicio:} Escribir de dos formas distintas la matriz:
$$
M=\left(
\begin{array}{lll}
  11 & 12 & 13 \\
  21 & 22 & 23
\end{array}
\right)
$$
\end{frame}

% ___________________________________________________________________________

\hypertarget{sub-ndices}{}
\pdfbookmark[0]{Subíndices}{sub-ndices}
\begin{frame}[containsverbatim]\frametitle{Subíndices}

  Los arrays son indexables del modo usual (uso de paréntesis, \texttt{1} es el primer índice)
\begin{Octave}
  >> M(1,2) 
  ans = 12
  >> M(end,end) 
  ans = 23
\end{Octave}

\small
\vfill\hfill ...por supuesto, \texttt{end} es el \emph{ultimo} índice.
\end{frame}

% ___________________________________________________________________________


\begin{frame}[containsverbatim]
\frametitle{El lenguaje de Matlab/Octave es muuuy flexible}
\begin{Octave}
  >> a = pi 
  a = 3.1416
  >> a(1) 
  ans = 3.1416
  >> a(1,1) 
  ans = 3.1416
  >> a(1,1,1) 
  ans = 3.1416
\end{Octave}
\begin{itemize}
\item {} 
Lo anterior devolvería un error en cualquier otro lenguaje de programación.

\item {} 
¿Esto es bueno o es malo?

\end{itemize}

\end{frame}

% ___________________________________________________________________________

\begin{frame}[containsverbatim]
\frametitle{Al menos, hay verificación de límites}
\begin{Octave}
  >> a=rand(3,3) 
  a = 
 
    0.583220   0.285967   0.010804 
    0.336715   0.530313   0.860391 
    0.060404   0.999098   0.501339 
 
    >> a(3,4) 
    error: invalid column index = 4
\end{Octave}
  
\end{frame}

% ___________________________________________________________________________

\hypertarget{pero-no-para-la-asignaci-n}{}
\pdfbookmark[0]{¡Pero no para la asignación!}{pero-no-para-la-asignaci-n}
\begin{frame}[containsverbatim]
\frametitle{¡Pero no para la asignación!}

\begin{Octave}
  >> nuevo(4)=1 
  nuevo = 
 
     0   0   0   1 
 
  >> nuevo(4,3)=2 
  nuevo = 
 
     0   0   0   1 
     0   0   0   0 
     0   0   0   0 
     0   0   2   0
\end{Octave}


\end{frame}

% ___________________________________________________________________________

\hypertarget{submatrices}{}
\pdfbookmark[0]{Submatrices}{submatrices}
\begin{frame}\frametitle{Submatrices}

Supongamos que de la matriz $M$ queremos extraer una
matriz $S$ formada por las 3 últimas filas y columnas:

$$
M=\left(\begin{array}{rrrrr}
  11 & 12 & 13 & 14 & 15 \\
  21 & 22 & 23 & 24 & 25 \\
  31 & 32 & \alert{33} & \alert{34} & \alert{35} \\
  41 & 42 & \alert{43} & \alert{44} & \alert{45} \\
  51 & 52 & \alert{53} & \alert{54} & \alert{55}
\end{array}\right)
$$
\end{frame}

% ___________________________________________________________________________

\begin{frame}[containsverbatim]
\frametitle{Submatrices}

Los subíndices se seleccionan mediante secuencias.  Por ejemplo, para
extraer la submatriz formada por las filas
\begin{itemize}
\item \emph{i} de 3 a 5
\item \emph{j} de 3 a 5
\end{itemize}
escribimos:
\vspace{-0.6em}
\begin{Octave}
>> S = M(3:5,3:5) 
 33   34   35 
 43   44   45 
 53   54   55
\end{Octave}

\medskip
Para extraer elementos de una fila o columna, podemos usar
'\verb|:|'. Por ejemplo:
\vspace{-0.5em}
\begin{Octave}
>> S = M(3:5,:) 
 31   32   33   34   35 
 41   42   43   44   45 
 51   52   53   54   55
\end{Octave}


\end{frame}

% ___________________________________________________________________________


\begin{frame}[containsverbatim]
\frametitle{Submatrices}

\textbf{Ejercicio}: La orden \verb|pascal(N)| devuelve una matriz
cuadrada de orden $N$ formada a partir del respectivo triángulo de
Pascal. 
\begin{itemize}
\item Calcular la matriz de Pascal $P$ de orden $6$.
\item 
Extraer la matriz $T$ formada por:
\begin{itemize}
\item Las filas $1$ a $3$
\item Las columnas $2$ y $5$
\end{itemize}
\end{itemize}
\end{frame}

\begin{frame}[containsverbatim]
\frametitle{Submatrices}

\textbf{Solución}: La matriz formada por las
\begin{itemize}
\item filas $1$ a $3$ y...
\item columnas $2$ y $5$
\end{itemize}
se obtiene de la siguiente forma:

\begin{Octave}
>> P = pascal(6)
>> T = P(1:3,[2,5]);
\end{Octave}

\end{frame}

% ___________________________________________________________________________

\begin{frame}[containsverbatim]
\frametitle{Submatrices}
\textbf{Ejercicio}: Ahora queremos extraer la matriz formada por las
columnas $1$, $3$ y $5$.
\hfill
\end{frame}

\begin{frame}[containsverbatim]
\frametitle{Submatrices}
\textbf{Ejercicio}: Ahora queremos extraer la matriz formada por las
columnas $1$, $3$ y $5$.

\hfill

\textbf{Solución}: La matriz formada por las
columnas $1$, $3$ y $5$ es...

\begin{Octave}
>> P(:,1:2:5) 
\end{Octave}
 
\end{frame}

% ___________________________________________________________________________

\hypertarget{operadores}{}
\pdfbookmark[0]{Operadores}{operadores}
\begin{frame}[containsverbatim]
\frametitle{Operadores}
\begin{description}[itemsep=1.1em]
\item [Operadores matriciales:]
\verb|+    -    *    /    ^|

\item [Operadores escalares:]
\verb|.*    ./    .^|

\item [Operadores lógicos:]
\verb# &   |    ! #

\item [Operadores de comparación:]
\verb|<    >    ==    >=    <=    !=|

\item [Operadores matriciales o de conjuntos:]
\verb# &&   ||#

\end{description}

\end{frame}

% ___________________________________________________________________________

\begin{frame}[containsverbatim]
\frametitle{Diferencias entre \& y \&\&}
\begin{Octave}
>> a=[1,2;0,1]; 
>> b=[1,0;0,1]; 
>> a&b 
ans = 
 
   1   0 
   0   1 
 
>> a&&b 
ans = 0
\end{Octave}

\begin{itemize}
\item 
  \verb|&&| se utiliza para las sentencias condicionales.
\item 
  \verb|&&| es un ``operador de cortocircuito''
\item 
  \textbf{Ejercicio}: ¿Cuándo es \verb'A && A' verdadero? 
%  ¿y \verb' A || A'? ¿y \verb'A && ~A'? ¿y \verb'A || ~A'?
  La orden \verb'false && 1/0' ¿devuelve un error (división por cero)?
\end{itemize}
\end{frame}

% ___________________________________________________________________________

\hypertarget{operadores-ii}{}
\pdfbookmark[0]{Operadores II}{operadores-ii}
\begin{frame}[containsverbatim]
\frametitle{Operadores II}

Operadores matriciales y escalares pueden confundirse fácilmente:
\begin{Octave}
>> a=rand(3,3);b=rand(3,3); 
>> a*b 
ans = 
    1.0297    0.9105    0.3293 
    0.9663    0.8267    0.4211 
    0.5355    0.4318    0.3279 
>> a.*b 
ans = 
    0.1824    0.3253    0.0563 
    0.5500    0.6003    0.1897 
    0.0458    0.0017    0.1822
\end{Octave}
  
\end{frame}

% ___________________________________________________________________________

\hypertarget{operadores-iii}{}
\pdfbookmark[0]{Operadores III}{operadores-iii}
\begin{frame}[containsverbatim]
\frametitle{¿Cuál es la solución ``correcta''?}

\begin{Octave}
>> a=[1,2,3;4,5,6;7,8,9]; 
>> a.^pi 
ans = 
    1.0000    8.8250   31.5443 
   77.8802  156.9925  278.3776 
  451.8079  687.2913  995.0416 
>> a^pi 
ans = 
 1.0e+03 * 
 0.6943-0.0004i 0.8540-0.0001i 1.0136+0.0002i 
 1.5743-0.0000i 1.9344-0.0000i 2.2946+0.0000i 
 2.4543+0.0003i 3.0149+0.0001i 3.5756-0.0002i
\end{Octave}

\end{frame}

% ___________________________________________________________________________

\begin{frame}\frametitle{Álgebra}

Con estas funciones se puede crear casi cualquier matriz
\begin{description}
\item [eye:]
Matriz de ceros con unos en la diagonal


\item [linspace:]
Vector de elementos equiespaciados


\item [logspace:]
Vector de elementos espaciados exponencialmente


\item [ones:]
Matriz de unos

\item [diag:]
Matriz con la diagonal que decidamos

\item [zeros:]
Matriz de ceros


\item [rand:]
Matriz de números pseudoaleatorios.


\item [meshgrid:]
Matrices equiespaciadas de dos dimensiones

\end{description}

\end{frame}

% ___________________________________________________________________________

\hypertarget{manipulaci-n-de-matrices}{}
\pdfbookmark[0]{Manipulación de matrices}{manipulaci-n-de-matrices}
\begin{frame}\frametitle{Álgebra}
Manipulación de matrices
\begin{description}
\item [reshape:]
Cambia la forma de la matriz conservando el número de
elementos


\item [transpose:]
Traspuesta. Equivale a .'


\item [ctranspose:]
Matriz conjugada. Equivale a '


\item [rot90:]
Gira la matriz 90 grados en sentido antihorario.


\end{description}
Tamaño de una matriz
\begin{description}
\item[size]  Devuelve el número de filas y columnas
\item[length] Devuelve el número de elementos de un vector\footnote{En
    realidad, \texttt{length(A)=max(size(A))}}
\end{description}

\end{frame}

% ___________________________________________________________________________

\begin{frame}[containsverbatim]
\frametitle{Resolución de sistemas lineales}

Para resolver sistemas de eucuaciones lineales contamos con un
operador universal:
\begin{Octave}
>> A=[1,0;2,1]; y=[2;4]; 
>> x=A\y 
x = 
  2 
  0
\end{Octave}
\begin{itemize}\item 
  A bajo nivel se usan las (excelentes) bibliotecas ATLAS y LAPACK
\end{itemize}
\end{frame}

%%% Local Variables: 
%%% mode: latex
%%% TeX-master: "modulo_octave"
%%% End: 


\section{Representación de gráficos}

\subsection{Gráficos 2D}

% ___________________________________________________________________________

\begin{frame}\frametitle{Representación gráfica}
\begin{itemize}
\item {} 
Repesentar datos gráficamente en Octave es sencillo e intuitivo.
\item  
  Se pueden usar distintos sistemas gráficos a través de un conjunto de funciones compatibles con Matlab.
\item Verremos un ejemplo en 2D
\end{itemize}

\end{frame}

% % ___________________________________________________________________________

\begin{frame}[containsverbatim]
\frametitle{Plot 2D}

La manera más sencilla de representar datos es mediante la función
\emph{plot}
\begin{description}
\item [Ejemplo:]
Representar una curva en el plano a partir de dos series de
datos:
\begin{Octave}
>> x=linspace(0,20,100); 
>> y=exp(-x./100).*sin(x);
>> plot(x,y);
\end{Octave}
\end{description}
\end{frame}

% ___________________________________________________________________________

\begin{frame}\frametitle{Plot 2D (II)}

El resultado (similar a)...
\begin{figure}[htbp]\begin{center}
%\includegraphics[width=600px]{files/abanico.jpg}
\pgfimage[width=0.8\linewidth]{img/plot-exp-sin.png}
\end{center}\end{figure}

\end{frame}

% ___________________________________________________________________________


\begin{frame}[containsverbatim]
\frametitle{Plot 2D (III)}

Los atributos de las gráficas se introducen con la ventana activa
\begin{description}
\item [Ejemplo:]

\begin{Octave}
>> title('Una funcion cualquiera') 
>> xlabel('Tiempo') 
>> ylabel('Amplitud')
\end{Octave}

\end{description}

\end{frame}


% ___________________________________________________________________________

\hypertarget{plot-iv}{}
\pdfbookmark[0]{Plot 2D (IV)}{plot-iv}
\begin{frame}\frametitle{Plot IV}

El resultado...
\begin{figure}[htbp]\begin{center}

\pgfimage[width=0.8\linewidth]{img/abanico2.png}
\end{center}\end{figure}

\end{frame}

% ___________________________________________________________________________

\hypertarget{plot-v}{}
\pdfbookmark[0]{Plot 2D (V)}{plot-v}
\begin{frame}[containsverbatim]
\frametitle{Plot V}
\small
Podemos poner varias curvas con distintos estilos:
\begin{Octave}
>> set(0, 'DefaultLineLineWidth', 1.5);   % Grosor 
>> x = linspace(0,20,50);
>> plot (x,exp(x/10).*sin(x),':',...
         x,exp(x/10).*cos(x),'-^'); 
\end{Octave}

Añadir una leyenda a posteriori...
\begin{Octave}
>> legend('exp(x/10)*sin(x)','exp(x/10)*cos(x)');
>> replot;  % Volver a pintar la gr'afica
\end{Octave}

Y exportar a un fichero
\begin{Octave}
  print -deps dibujo.eps
\end{Octave}

En la ayuda de \verb|print| está la lista de dispositivos (\verb|-d|)
\end{frame}

% ___________________________________________________________________________

\begin{frame}\frametitle{Plot 2D (VI)}
  El resultado...
\begin{center}
\pgfimage[width=0.85\linewidth]{img/plot-exp-sin-cos.png}
\end{center}
\end{frame}

\subsection{Gráficos en 3D}


\begin{frame}[containsverbatim]
  \frametitle{Curvas 3D}
  Las gráficas 3D más sencillas\footnote{Aunque
  \texttt{plot3} tiene otras posibilidades, por ejemplo representar
  familias de curvas (cuando \texttt{x}, \texttt{y}, \texttt{z} son
  matrices). Véase la ayuda.} son curvas trazadas mediante la función
  \emph{plot3}
  \begin{description}
  \item [Ejemplo:] Representar una curva en el espacio a partir de tres
    vectores de datos, \verb|x|, \verb|y|, \verb|z|.
\begin{Octave}
  >> t=0:0.1:10*pi;
  >> x=(1+sin(t)).*cos(t);
  >> y=(1+sin(t)).*sin(t);
  >> z=t;
  >> plot3(x,y,z);
\end{Octave}
\end{description}
\end{frame}

\begin{frame}\frametitle{Curvas 3D: El resultado...}
\begin{center}
\pgfimage[width=0.9\linewidth]{img/plot3.png}
\end{center}
\end{frame}

\begin{frame}[containsverbatim]
  \frametitle{Superficies 3D: Generación de mallas}
  La función \verb|[x,y]=meshgrid(xmin:hx:xmax, ymin:hy:ymax)| genera
  dos matrices \verb|x| e \verb|y|:
  \begin{itemize}
  \item Filas de \verb|x|: copias del vector \verb|xmin:hx:xmax|
  \item Columnas de \verb|y|: vector \verb|(ymin:hy:ymax)'|
  \end{itemize}
  Ejemplo
\begin{Octave}
  >> [x,y]=meshgrid(0:2,0:1)
  x =
   0   1   2
   0   1   2

 y =
   0   0   0
   1   1   1
\end{Octave}
Los \verb|(x[i],y[j])| son un mallado de \verb|[xmin,xmax]|$\times$\verb|[ymin,ymax]|
\end{frame}


\begin{frame}[containsverbatim]
  \frametitle{Dibujo de superficies 3D}
  La orden \verb|mesh(x,y,z)| representa una superficie 3D, siendo
  \begin{itemize}
  \item \verb|x|, \verb|y| dos matrices que representan una malla
  \item \verb|z| otra matriz que representa los valores $f(x,y)$
  \end{itemize}
  Ejemplo: gráfica de la función $f(x,y)=sin(x)*cos(y)$
  \begin{Octave}
    >> [x,y]=meshgrid(-pi:0.1:pi, -pi:0.1:pi);
    >> z=sin(x).*sin(y);
    >> mesh(x,y,z);
  \end{Octave}
  Podemos cambiar el punto de vista con la función \verb|view(a,e)|,
  donde 
  \begin{itemize}
  \item \verb|a| es el \emph{azimut} (ángulo de rotación en el plano xy)
  \item \verb|e| es la \emph{elevación} vertical
  \end{itemize}
  \begin{Octave}
  >> view(30,60);
  >> colorbar('East'); % Pintar una barra vertical
  \end{Octave}
\end{frame}

\begin{frame}\frametitle{Superficies 3D: El resultado...}
\begin{center}
\pgfimage[width=0.9\linewidth]{img/mesh-3d.png}
\end{center}
\end{frame}

\begin{frame}[containsverbatim]
  \frametitle{Curvas de nivel}
  \begin{itemize}
  \item
    La orden \verb|contour(x,y,z)| representa curvas de nivel
  \item \verb|x|, \verb|y| y \verb|z| tienen el mismo significado que
    en la orden \verb|mesh|
  \item Acepta un cuarto parámetro, \verb|n| (número de curvas)
  \item La variante \verb|contourf| dibuja un degradado de colores
  \end{itemize}
  \begin{Octave}
  >> [x,y]=meshgrid(-pi:0.1:pi, -pi:0.1:pi);
  >> z=sin(x).*sin(y);
  >> contourf(x,y,z);    
  \end{Octave}
\end{frame}

\begin{frame}\frametitle{Curvas de nivel: El resultado...}
\begin{center}
\pgfimage[width=0.9\linewidth]{img/contour.png}
\end{center}
\end{frame}


% \begin{frame}[containsverbatim]
% \frametitle{Mallas estructuradas 3D con VTK}
% \begin{enumerate}
% \item Creamos un array 3D y lo grabamos en formato \verb|hdf5|
%   \begin{Octave}
% nx=ny=nz=4; % Por ejemplo, 4 ptos en cada eje
% A=rand(nx,ny,nz); % array aleatorio nx*ny*nz
% save -hdf5 "malla_estruc.h5" A;
%   \end{Octave}
% \item Convertimos el fichero \verb|hdf5| en formato \verb|vtk|
%   \begin{sh}
% $ h5tovtk -d /A/value malla_estruc.h5
%   \end{sh}  % $fin de modo matemático, editor!
%   \begin{itemize}
%   \item \footnotesize Se ha usado el programa \texttt{h5tovtk} (en
%     Ubuntu, paquete h5utils).
%   \item \footnotesize La opción ``\texttt{-d}'' indica la ruta hacia
%     el ``dataset'', dentro del fichero (ésta se puede localizar con
%     el programa \texttt{h5dump} o \texttt{h5ls -r})
%   \end{itemize}
% \item Visualizamos el fichero resultante (\texttt{malla\_estruc.vtk})
%   con \texttt{mayavi2} o con \texttt{paraview} (filtro ``\emph{Treshold}'')
% \end{enumerate}
% \end{frame}
%
% \begin{frame}\frametitle{Resultado (generado con paraview)}
% \begin{center}
% \pgfimage[width=0.9\linewidth]{img/cubo.png}
% \end{center}
% \end{frame}


%%% Local Variables: 
%%% mode: latex
%%% TeX-master: "modulo_octave"
%%% End: 


\section{Programación con Octave}
\hypertarget{control-de-flujo}{}
\pdfbookmark[0]{Control de flujo}{control-de-flujo}
\begin{frame}\frametitle{Control de flujo}
\begin{itemize}
\item {} 
No hay diferencias con otros lenguajes de programación
\emph{interpretados}
\item {} 
En estos lenguajes un bucle es en realidad un \emph{iterador}
\end{itemize}

\end{frame}

% ___________________________________________________________________________

\hypertarget{condicionales}{}
\pdfbookmark[0]{Condicionales}{condicionales}
\begin{frame}[containsverbatim]
\frametitle{Condicionales}

Este es un ejemplo del uso de la sentencia \emph{if}:
\begin{Octave}
if saludo 
  disp('hola') 
else 
  disp('no te saludo') 
end
\end{Octave}

\begin{itemize}
\item 
  ¿Cuál es la salida si saludo=1?
\item 
  En Octave, podemos usar \texttt{endif} (no compatible Matlab!).
  Igual ocurrirá con \texttt{endfor}, \texttt{endfunction}
\end{itemize}
\end{frame}

% ___________________________________________________________________________

\hypertarget{bucles}{}
\pdfbookmark[0]{Bucles}{bucles}
\begin{frame}[containsverbatim]
\frametitle{Bucles}

Ejemplo del uso de la sentencia \emph{for}:
\begin{Octave}
function primetest(n) 
  printf('Numeros primos de 1 a %i\n',n) 
  for i=1:n 
    if isprime(i) 
      disp(i) 
    endif
  endfor
endfunction
\end{Octave}
\textbf{Ejercicio}:
\begin{itemize}
\item Programar la función en el fichero
  ``\texttt{primetest.m}''
\item 
  ¿Cuál es la salida de esta subrutina si n=20?
\end{itemize}
\end{frame}

% ___________________________________________________________________________

\hypertarget{bucles-iteradores}{}
\pdfbookmark[0]{Bucles iteradores}{bucles-iteradores}
\begin{frame}[containsverbatim]
\frametitle{Bucles iteradores}

Recordad que secuencia == vector
\begin{Octave}
>> for i=[1,3,2,4,3,4,5] 
> disp(i) 
> end 
 1 
 3 
 2 
 4 
 3 
 4 
 5
\end{Octave}

\end{frame}

% ___________________________________________________________________________

\hypertarget{otras-sentencias}{}
\pdfbookmark[0]{Otras sentencias}{otras-sentencias}
\begin{frame}\frametitle{Otras sentencias}
\begin{quote}
\begin{description}

\item [while:]
Bucle controlado por condición lógica

\item [case:]
Control de casos cerrados

\item [try:]
Sentencia de control para probar errores

\item [break:]
Permite abandonar directamente el interior de un bucle


\item [continue:]
Se utiliza para saltar al final de la actual iteración de un bucle


\item [return:]
Devuelve el control al programa principal

\end{description}
\vfill
\textbf{Ejercicio} Realizar distintos ejemplos con las
sentencias anteriores
\end{quote}

\end{frame}



% ___________________________________________________________________________

\begin{frame}[fragile]{Funciones en MATLAB}


\lstset{
  language=Matlab,
  basicstyle=\ttfamily\footnotesize,
  keywordstyle=\color{blue!70!black}\bfseries,
  commentstyle=\color{green!40!black}\itshape,
  frame=single,
  framerule=0.3pt,
  columns=fullflexible,
  keepspaces=true,
  showstringspaces=false,
  numbers=none
}

\begin{columns}[T]
  \begin{column}{0.52\textwidth}
    \textbf{1. Función en fichero} (\texttt{\small area\_circulo.m})
\begin{lstlisting}
function [A, P] = area_circulo(r)
% AREA_CIRCULO  Area y perimetro
%   [A, P] = area_circulo(r)
  A = pi * r.^2;
  P = 2 * pi * r;
end
\end{lstlisting}
    \vspace{0.4em}
    \textbf{Uso:}\vspace{-0.5em}
\begin{lstlisting}
[A, P] = area_circulo(2);
A = area_circulo(1:3);
\end{lstlisting}
  \end{column}

  \begin{column}{0.46\textwidth}
    \textbf{2. Función anónima}
\begin{lstlisting}
f = @(x) x.^2 + 1;
f(3)       % 10
a = 2;
g = @(x) a*x;  % captura a
\end{lstlisting}
    \vspace{0.4em}
    \textbf{Como argumento:}
\begin{lstlisting}
fzero(@(x) cos(x)-x, 0.5)
integral(@(x) exp(-x.^2), 0, 1)
\end{lstlisting}
  \end{column}
\end{columns}

\vspace{0.6em}
\begin{itemize}
  \item El fichero debe llamarse \texttt{area\_circulo.m}: manda el nombre del fichero.
  \item Varias salidas entre corchetes; \texttt{nargin} y \texttt{nargout} dan el n.º de argumentos.
  \item Usa \texttt{.*}, \texttt{./}, \texttt{.\^{}} para que la función admita vectores.
  \item Una anónima es una sola expresión y congela las variables externas al definirla.
\end{itemize}
\end{frame}


% ___________________________________________________________________________

\hypertarget{funciones}{}
\pdfbookmark[0]{Funciones}{funciones}
\begin{frame}[containsverbatim]
  \frametitle{Las funciones con algo más de profundidad}

Varios valores devueltos
\begin{Octave}
function [m,M] = minmax(valores) 
  % Esta funci'on toma como par'ametro un vector y 
  % devuelve sus valores m'inimo y m'aximo
  m = min(valores); 
  M = max(valores);
endfunction 
 
help minmax
lista = 1./(1:10);
m = minmax(lista);
\end{Octave}

\textbf{Observaciones}
\begin{itemize}
\item Recordar que \texttt{endfunction} no es compatible con Matlab
\item \texttt{endfunction} es opcional si función está en su fichero 
\end{itemize}

\end{frame}

\begin{frame}[fragile]
\frametitle{Ejercicio}
\begin{itemize}
    \item Programa una función llamada \texttt{integracion\_numer(f, a, b, N)}, que aplique la fórmula de los trapecios compuesta\footnote{\url{https://es.wikipedia.org/wiki/Regla_del_trapecio\#Regla_del_trapecio_compuesta}} con \texttt{N} subintervalos para aproximar la integral de una función $f$ en un intervalo $(a,b)$.
  \item Calcula de forma manual la integral de algunas funciones sencillas, por ejemplo $f_1(x)=x^2$, $f_2(x)=\cos(x)$, $f_3(x)=3\cos(3x)\cdot e^{\sin(3x)}$.
  \item Aproxima estas integrales con la función que hemos definido con diferentes valores de $N$ y evalúa el error respecto a la interal exacta que calculaste
  \item Genera una tabla con los resultados
\end{itemize}
\end{frame}

% ___________________________________________________________________________

\hypertarget{funciones-an-nimas}{}
\pdfbookmark[0]{Funciones anónimas.}{funciones-an-nimas}
\begin{frame}[containsverbatim]
\frametitle{Más sobre funciones anónimas}
\small
Definición
\begin{itemize}
\item Matlab como casi todos los senguajes, define un tipo especial de funciones,
  (\emph{funciones lambda}). 
 Para ello, se utiliza el carácter \verb|@|
\item Permite definr funciones breves sin necesidad de un archivo
  adicional\footnote{Matlab no permite la definición de funciones de
    forma interactiva (salvo \texttt{inline}). Octave sí lo permite.}

\end{itemize}
\begin{Octave}
>> testf = @(x,y) exp(-(x.^2+y.^2)) 
>> testf(1,i) 
ans = 1
\end{Octave}
Utilidad: estas funciones  pueden pasarse como argumento para otra
función
\begin{Octave}
>> integracion_numer(@(x) x^2, 0, 1);
\end{Octave}
\end{frame}


% ___________________________________________________________________________

\begin{frame}[containsverbatim]
\frametitle{Problemas arreglados por Octave}

Declaración de funciones directamente en consola:
\begin{Octave}
>> function y=foo(x) 
> y=3*x 
> endfunction
>> foo(4) 
y =  12 
ans =  12
\end{Octave}
\end{frame}

\begin{frame}[containsverbatim]
 \frametitle{Problemas arreglados por Octave}
 Varias funciones en un mismo archivo.
 \begin{itemize}
 \item 
 Definimos el fichero \texttt{biblioteca.m}
 \begin{Octave}
   function y=cuadrado(x)
     y=x^2;
   end
   function y=raiz(x)
     y=sqrt(x);
   end
 \end{Octave}
 \item 
 Podemos acceder a las dos funciones:
 \begin{Octave}
   >> source('biblioteca.m')
   >> cuadrado(raiz(2))
   ans =  2.0000
 \end{Octave}
\end{itemize}
\end{frame}



% ___________________________________________________________________________
\begin{frame}[containsverbatim]
\frametitle{Problemas arreglados por Octave}

Ejecución de ``funciones compuestas''
\begin{Octave}
>> a=rand(3,3)(1,2) 
a =  0.14272
\end{Octave}

Permite expresiones mucho más compactas.

\end{frame}


% ___________________________________________________________________________

\hypertarget{la-biblioteca-de-funciones}{}
\pdfbookmark[0]{La biblioteca de funciones}{la-biblioteca-de-funciones}
\begin{frame}\frametitle{La biblioteca de funciones}
\begin{itemize}
  \setlength{\itemsep}{0.7em}
\item {} 
  Integración numérica (\textbf{quad}, quadl, quad2dg...)

\item {} 
  Interpolación (\textbf{interp1}, interp2, polyfit, fft,...)

\item {} 
  Resolución numérica de ecuaciones diferenciales (\textbf{ode45},...)

\item {} 
Estadística descriptiva (mean, std, median,...)

\end{itemize}

% \textbf{Ejercicio}: Localiza funciones de tu interés y pruébalas!
% \begin{itemize}\item 
%   Indicación:
%   \begin{itemize}
%   \item Echa un vistazo a la ayuda de Octave (HTML)
%   \item Utiliza la función \texttt{lookfor} (busca funciones relacionadas
%     con una palabra clave). Por ejemplo, \texttt{lookfor fourier}.
%   \end{itemize}

% \end{itemize}


\end{frame}

% ___________________________________________________________________________

% \hypertarget{toolkits-en-octave}{}
% \pdfbookmark[0]{``Toolkits'' en Octave}{toolkits-en-octave}
% \begin{frame}\frametitle{``Toolkits'' en Octave}
% \begin{itemize}
% \item {} 
% \texttt{Tookits} ``oficiales'': Octave forge \href{http://octave.sourceforge.net}{http://octave.sourceforge.net}
% \begin{itemize}
% \item Mirar en la web la lista de \texttt{toolkits}
% \end{itemize}
% \item {} 
% Muchos otros.
% \begin{quote}
% \begin{itemize}
%
% \item MPITB (\texttt{Toolkit} de Paralelización con MPI, desarrollado por
%   Javier F. Baldomero, Universidad de Granada)\footnote{\url{http://atc.ugr.es/javier-bin/mpitb}}
%
% \item {} 
% Octaviz (VTK en Octave)
%
% \item {} 
% Octplot, etc
%
% \end{itemize}
% \end{quote}
%
% \item Problema: no hay garantías de compatibilidad con \texttt{Toolkits} de Matlab
%   \begin{itemize}
%   \item Independencia respecto a Octave y sus líneas estratégicas
%   \end{itemize}
% \end{itemize}
%
% \end{frame}
%
% % ___________________________________________________________________________

% \hypertarget{ejercicios}{}
% \pdfbookmark[0]{Ejercicios}{ejercicios}
% \begin{frame}\frametitle{Ejercicios}

% \end{frame}

% % ___________________________________________________________________________

% \hypertarget{ejercicio}{}
% \pdfbookmark[0]{Ejercicio}{ejercicio}
% \begin{frame}\frametitle{Ejercicio}

% Resolver la ecuación de Laplace en 1d con CC Dirichlet:

% Mediante diferencias finitas.

% \end{frame}

% % ___________________________________________________________________________

% \hypertarget{id6}{}
% \pdfbookmark[0]{Ejercicio}{id6}
% \begin{frame}\frametitle{Ejercicio}

% \end{frame}

% % ___________________________________________________________________________

% \hypertarget{id7}{}
% \pdfbookmark[0]{Ejercicio}{id7}
% \begin{frame}\frametitle{Ejercicio}

% Crear una \emph{función} que genere esta matriz con \emph{n} arbitrario

% Usando también la función \emph{diag}

% \end{frame}

% % ___________________________________________________________________________

% \hypertarget{id8}{}
% \pdfbookmark[0]{Ejercicio}{id8}
% \begin{frame}\frametitle{Ejercicio}

% Crear una función que cree el término inEependiente con \emph{n}
% arbitrario:

% y resolver para \emph{n} = 7

% \end{frame}

% ___________________________________________________________________________

% \begin{frame}[containsverbatim]
% \frametitle{Extensiones dinámicas de Octave}
% \vspace{-0.5em}
%   Fichero \texttt{MiFuncion.cpp}:
%
%   \begin{C++}
% #include <oct.h> // Acceso al codigo de Octave!
% DEFUN_DLD(MiFuncion, args, nargout, "Descripcion")  
% {
%   // ... Codigo C++ ...
%   return octave_value()
% }
%   \end{C++}
%
%   Compilación (con \texttt{mkoctfile}, un envoltorio del compilador  C++)
%
%   \begin{sh}
% $ mkoctfile MiFuncion.cpp
%   \end{sh} %$
%
%   Obtenemos el fichero \texttt{MiFuncion.oct}, que usaremos desde octave:
%   \begin{Octave}
% >> MiFuncion
%   \end{Octave}
% \end{frame}
%
%
% \begin{frame}[containsverbatim]
% \frametitle{Extensiones dinámicas de Octave: Hola Mundo!}
%
% Fichero \texttt{holamundo.cpp}
% \begin{C++}
% #include <octave/oct.h>
% DEFUN_DLD(holamundo, args, nargs , ``Programa hola mundo'')
% {
%   std::cout << ``Hola Mundo!'' << std::endl;
%   return octave_value()
% }
% \end{C++} 
%
% Lo compilamos (\verb|$ mkoctfile holamundo.cpp|) y lo usamos desde Octave:
%   \begin{Octave}
% >> holamundo
% ``Hola Mundo!''
%   \end{Octave}
% \end{frame}
%
% \begin{frame}[containsverbatim]
% \frametitle{Extensiones dinámicas de Octave: $f(x)=x^2$ }
% \small
% Fichero \texttt{xcuadrado.cpp}:
% \begin{C++}
% #include <octave/oct.h>
% DEFUN_DLD(xcuadrado, args, , \
%   "Calcula el cuadrado de un real, x'')
% {
%   if(args.length() != 1) {
%     std::cerr << "ERROR: Se necesita tomar un valor" << std::endl;
%     return octave_value()
%   }
%   const double x = args(0).double_value();
%   octave_value y = x*x;
%   return y;
% }
% \end{C++}
% \end{frame}
%
% \begin{frame}[containsverbatim]
%   \frametitle{Creación de módulos dinámicos C/C++ con SWIG} 
%
%   \begin{itemize}
%   \item 
%   \textbf{¿Qué es SWIG?} 
%   Potente interfaz para generar envolturas (\textit{wrappers}) que
%   permiten acceder a funciones (o clases) C/C++ desde Octave, Python,
%   Java, Perl,...
%   \item \textbf{¿Cómo usar SWIG?}
%     \begin{enumerate}
%     \item Tenemos varias funciones (o clases) C/C++ en un fichero
%       ``\verb|ejemplo.cpp|''
%     \item Escribimos un fichero de directivas ``\verb|ejemplo.i|'', donde
%       indicamos a qué funciones C/C++ deseamos acceder
%     \item A partir de él, SWIG genera un fichero C/C++ ``\verb|ejemplo_wrap.cpp|''
%     \item Compilamos ``\verb|ejemplo_wrap.cpp|'' $+$
%       ``\verb|ejemplo.cpp|'' (mediante \verb|mkoctfile|), obteniendo
%       ``\verb|ejemplo.oct|''
%     \item Desde Octave, ya podemos utilizar el módulo o \emph{toolkit}
%       ``\verb|ejemplo|'' 
%     \end{enumerate}
%   \end{itemize}
% \end{frame}
% ___________________________________________________________________________

\begin{frame}\frametitle{Bibliografía}
\begin{itemize}
\item {}
\textit{Introducción informal a Matlab y Octave}:\par
\href{http://iimyo.forja.rediris.es/}{http://iimyo.forja.rediris.es/}

\item {}
\emph{A. Quarteroni, F. Salieri} Cálculo científico con Matlab y
Octave. Springer 2006.

\item {}
\textit{Documentación ``oficial'':}\par
\href{http://www.gnu.org/software/octave/docs.html}{http://www.gnu.org/software/octave/docs.html}

\item {} 
\textit{Documentación de Octave-Forge:}\par
\href{http://octave.sourceforge.net/docs.html}{http://octave.sourceforge.net/docs.html}
\begin{itemize}
\item {} 
Incluye
\href{http://octave.sourceforge.net/coda/coda.pdf}{http://octave.sourceforge.net/coda/coda.pdf},
manual para extensiones dinámicas en C++.

\end{itemize}

\end{itemize}
\end{frame}

%%% Local Variables: 
%%% mode: latex
%%% TeX-master: "modulo_octave"
%%% End: 


\end{document}

%%% Local Variables: 
%%% mode: latex
%%% TeX-master: t
%%% End: 
